\section{Hasil Eksperimen}
Tabel~\ref{tab:hasil} menyajikan hasil evaluasi untuk setiap konfigurasi
jumlah faktor laten pada model ALS.

\begin{table}[htb]
\centering
\caption{Hasil evaluasi model ALS pada berbagai jumlah faktor laten}
\label{tab:hasil}
\rowcolors{2}{uibg}{white}
\begin{tabular}{@{}crrr@{}}
\toprule
\textbf{Faktor Laten ($k$)} & \textbf{RMSE} & \textbf{Precision@10} & \textbf{NDCG@10} \\
\midrule
16  & 0,811 & 0,329 & 0,388 \\
32  & 0,768 & 0,357 & 0,419 \\
64  & 0,742 & 0,381 & 0,447 \\
128 & 0,745 & 0,376 & 0,441 \\
\bottomrule
\end{tabular}
\end{table}

Konfigurasi dengan 64 faktor laten memberikan performa terbaik pada ketiga
metrik. Peningkatan lebih lanjut ke 128 faktor tidak memberikan perbaikan
performa dan justru mengindikasikan gejala \emph{overfitting} ringan.
